\scanpage{105}
\mspage{98}

\deleted{On va construire \deleted{deux} \inserted{\unsure{sous|deux}} systèmes inductifs \inserted{de $(X_i)$, $(Y_j)$} indexés par les ordinaux $\alpha<\alpha_0$, en procédant par récurrence transfinie.}

Ceci va résulter du

\uline{Lemme} 3 Soient $X=\varinjlim_I X_i$, $X\simeq\varinjlim_J Y_j$ deux représentations de $X$ sous forme de $\varinjlim$ filtrantes, avec $I$ grand devant $\pi'$, $J$ grand devant $\pi$, avec $\pi,\pi'$ deux cardinaux tels que $\pi'>\pi$. On suppose $(X_i)_{i\in I}$ $\pi'$-adapté, i.e. pour toute partie filtrante $I'$ de $I$ de cardinal $\leq\pi'$, $I'$ a une borne supérieure $i_{I'}$ et $X_{i_{I'}}=\varinjlim_{i\in I'} X_i$.

Alors pour \deleted{\unsure{tout}} $i\in I$, on peut trouver une partie $I'$ de $I$, \inserted{contenant $i$,} \uline{grande devant $\pi$}, de cardinal $\leq\pi'$, \deleted{bien ordonnée} et une partie $J'$ de $J$, \deleted{\unsure{munies} \unsure{d'une} application} \inserted{\unsure{munies|puis}} un isomorphisme $I'\xrightarrow{\overset{\unsure{\lambda|\varphi}}{\sim}}J'$ et une application \inserted{str.} croissante $\unsure{\mu|\psi}:J'\longrightarrow I'$, \deleted{tels que} et des hom

% archive sheet number: 104
